\scp{\tsc{Central Timor} internal diversity}{05-01}

Each of the \tsc{Central Timor} “languages” are internally diverse
with multiple dialects and/or varieties.
Kemak speakers report about at least a dozen different named
dialects of their language including: Atsabe,
Diirbatin, Kailaku, Kutubaba, Lemia,
Leolima, Leosibe, Marobo, Oho, and Saneri.
Data for about ten of these varieties is available,
with lexical similarity between any two varieties ranging from
about 76{\%} to 88{\%}.\footnote{
		Kutubaba Kemak is
		the most divergent known variety of Kemak.
		Kutubaba Kemak is spoken in Kenebibi village on the north coast
		of the easternmost Kemak-speaking area in \frf{05-01}.}

\citet[i]{LSG2006} lists three different varieties of Tokodede:
Bazartete Tokodede, Liquiçá Tokodede, and Maubara Tokodede.
Lexical differences are found between all three varieties as
well as phonological differences,
the most significant of which is Maubara Tokodede /v/ corresponding
to other Tokodede /h/.
One example is *babuy > Maubara Tokodede \it{vavi}, other Tokodede \it{hahi} ‘pig’.

\citet{Fogaca2017} shows that varieties of \tsc{Mambae} fall
into three groupings on the basis of lexical similarity: Northwest
Mambae, Central Mambae, and South Mambae.
The South Mambae grouping is most divergent with lexical similarity
compared to other varieties of \tsc{Mambae} ranging from 63{\%}--70{\%}.\footnote{
		The degree of lexical similarity between Northwest Mambae
		and South Mambae is similar to the degree of lexical similarity
		between Northwest Mambae and Tokodede, which is about 61{\%} \citep[35]{Edwards2019}.}
\citeauthor{Fogaca2017}’s figures summarising the degree of lexical similarity
between different varieties of \tsc{Mambae} are given in \trf{05-02}.


\begin{table}
\caption[\tsc{Mambae} lexical similarity (Fogaça 2017: 82)]
{\tsc{Mambae} lexical similarity \citep[82]{Fogaca2017}\su{†}}\label{tab:05-02}
%\begin{adjustbox}{width=\textwidth}
	\begin{threeparttable} %\st{2pt}
	\begin{tabular}{llllllllllll}
\hhline{---------}
\mc{3}{|l}{Bazartete}	&		&		&		&		&		&	\mc{1}{c|}{}	&&				 						\\
\mc{1}{|l}{81}	&	\mc{3}{l}{Hatulia}				&	\mc{5}{r|}{\tsc{Northwest Mambae}}	&	\\
\mc{1}{|l}{78}	&	81	&	\mc{2}{l}{Railaco}	&		&		&		&		&	\mc{1}{c|}{}	&&	\\ \hhline{-----------}
74	&	75	&	76	&	\mc{3}{|l}{Laulara}	&		&	& & &	\mc{1}{c|}{} \\%&		&		&		&	\\
74	&	76	&	76	&	\mc{1}{|l}{85}	&	\mc{3}{l}{Aileu Vila}	&		&		\mc{3}{r|}{\multirow{2}{*}{\tsc{Central Mambae}}}	\\
71	&	73	&	70	&	\mc{1}{|l}{81}	&	83	&	\mc{5}{l}{Hatu-Builico}	&			\mc{1}{c|}{}\\
69	&	69	&	70	&	\mc{1}{|l}{82}	&	80	&	82	&	\mc{4}{l}{Liquidoe}	&		\mc{1}{c|}{} \\ \hhline{~~~---------}
63	&	65	&	65	&	68	&	70	&	68	&	66	&	\mc{3}{|l}{Hatu-Udo}	&		&	\mc{1}{c|}{}\\
63	&	64	&	64	&	67	&	68	&	68	&	67	&	\mc{1}{|l}{82}	&	\mc{2}{l}{Letefoho}	&	\mc{2}{r|}{\tsc{South Mambae}}\\
63	&	66	&	63	&	66	&	68	&	67	&	65	&	\mc{1}{|l}{82}	&	80	&	Betano	& & \mc{1}{c|}{}\\ \hhline{~~~~~~~-----}
		\lspbottomrule
	\end{tabular}
		\begin{tablenotes}\tnsize
			\item[†] {Names of different varieties are from subdistricts where they are spoken.
								Letefoho and Betano are \it{sucos} (“villages”) in the subdistrict of Same.}
			%\item[‡]
			%\item[Δ]
			%\item[‖]
			%\item[¶]
		\end{tablenotes}
	\end{threeparttable}
%\end{adjustbox}
\end{table}

Because this chapter is focused on the internal relationships
of the languages of \tsc{Central Timor},
it necessarily highlights the similarities between these languages.
To give a sense of the structural,
phonological, and lexical diversity which is encompassed under
the label “\tsc{Mambae}” numerals in five different varieties
of \tsc{Mambae} are given in \trf{05-03}.

Structural diversity can be seen in the way the numerals 6--9
and multiples of 10 are formed.
While both Northwest Mambae
and South Mambae have complex forms for the numerals 6--9 (see
\srf{02-03-01-04}), they use different constructions.
Northwest Mambae uses \it{hoho} (no known independent meaning
--- also used in Tokodede for numerals 6--9).\footnote{
		Maubara Tokodede uses \it{vovo} for numerals 6--9 pointing to earlier **bobo.}
while South Mambae uses \it{liim nai} ‘five \it{nai}’.
Central Mambae differs again with mono-morphemic forms which
are ultimately from PMP.
For the multiples of ten, three different constructions are found:
Northwest Mambae \it{guul} X,
Central Mambae X \it{nulo} (or metathesised \it{nuul}),
and South Mambae \it{saguul haet} X ‘ten times X’ or just \it{haet} ‘times’.
The numerals in \trf{05-03} also give a sample of the phonological
diversity within the \tsc{Mambae} cluster.
Northwest Mambae /p/ corresponds to /f/ in Central Mambae
and either /p/ or /f/ in different varieties of South Mambae.
Similarly, Central Mambae has /k/ where other varieties have /g/.

\begin{table}
	\centering\tcp{Numerals in five varieties of \tsc{Mambae}}{05-03}
	\st{3.9pt}\begin{tabular}{llllll}\lsptoprule
	&	NW. Mambae	&	C. Mambae	&	C. Mambae	&	S. Mambae	&	S. Mambae\\
	&	Railaco	&	Liquidoe	&	Hatu-Builico	&	Hatu-Udo	&	Betano\\ \midrule
1	&	iid	&	iid	&	iid	&	iid	&	iid\\
2	&	ruu	&	ruu	&	ruu	&	rua	&	ruu\\
3	&	teul	&	teul	&	teul	&	teul	&	teul\\
4	&	paat	&	faat	&	faat	&	paat	&	faat\\
5	&	liim	&	liim	&	liim	&	liim	&	liim\\
6	&	hohon iid	&	neen	&	neen	&	liim nai niid	&	liim nai ida\\
7	&	hoho ruu	&	hitu	&	hitu	&	liim nai rua	&	liim nai rua\\
8	&	hoho teul	&	ualu	&	ualu	&	liim nai telu	&	liim nai telu\\
9	&	hoho paat	&	sia	&	sian	&	liim nai pata	&	liim nai fata\\
10	&	saguul	&	sakuul	&	sakuul	&	saguul	&	saguul\\
20	&	guul ruu	&	rua nuul	&	rua nulo	&	saguul haet rua	&	haet rua\\
30	&	guul teul	&	teul nuul	&	teul nulo	&	saguul haet teul	&	haet teul\\
100	&	atus iid	&	atus iid	&	atus iid	&	atus iid	&	atus iid\\
		\lspbottomrule
	\end{tabular}
\end{table}

Further diversity within \tsc{Mambae} can be seen by examining
the data in \citet{Fogaca2017}.
Northwest Mambae and Central Mambae retain a genitive suffix
\it{-n} which has been lost in South Mambae; e.g.
Northwest/Central Mambae \it{ate-n}, South Mambae \it{ate} ‘liver’.
Similarly, Northwest Mambae and Central Mambae have a nominal
suffix/enclitic \it{-a} of unclear function; e.g.
Northwest/Central Mambae \it{ai-a}, South Mambae \it{ai} ‘tree’.
\footnote{
		Parallels between Northwest/Central Mambae \it{-a}
		and the \tsc{Rote-Meto} determiner **-a are striking.
		One function of this determiner in \tsc{Rote-Meto} languages
		is to mark definiteness.}

Another point of diversity is seen in the forms of metathesis.
All varieties of \tsc{Mambae} have productive CV → VC metathesis (e.g.
\it{lelo → leol} ‘sun’),
but different processes of  vowel assimilation operate after metathesis.
Thus, for instance, South Mambae from Betano village has lowering
of \it{i → e} after initial \it{a} when metathesis operates,
as seen in *hapuy > *api > \it{aep} ‘fire’.
The varieties of Northwest Mambae from Liquiçá
and Hatulia have taken this further with complete assimilation
of the initial vowel: *api > \it{eep-a} [ɛpa] ‘fire’.
Central Mambae, on the other hand,
does not appear to have any vowel assimilation in this particular
form with all known varieties having \it{aif-a} ‘fire’.

Throughout the remainder of the chapter,
we usually cite examples from only one variety each of Welaun,
Kemak, Tokodede, Northwest Mambae, Central Mambae, and South Mambae.
Unless otherwise stated, Welaun data is from the hamlet of Oele{\Q}u
in the village (\it{desa}) of Bauho /bauhoo/.
Kemak data is from the Kailaku variety,
as spoken in the village of Haikesak/Maumutin.
Tokodede data is from the Liquiçá /likisaa/ variety, which is
the focus of \citet{LSG2006}.
Northwest Mambae data are from the subdistrict of Railaco,
Central Mambae from Liquidoe,
and South Mambae from the village of Betano in the subdistrict of Same.
